\documentclass[11pt,a4paper]{article}
\usepackage{exercise-common}

\begin{document}

\exerciseTitle{Titulo del ejercicio}
\exerciseMeta{2bach}{Electromagnetismo}{Induccion electromagnetica}{Variacion de flujo magnetico}{3}{12}

\sectionLabel{Enunciado}

Este archivo es una plantilla para futuros enunciados. Sustituir este texto por un enunciado breve, claro y compacto. El PDF debe parecerse a un enunciado de PAU/selectividad: poco espacio decorativo, figura cerca del texto y apartados visibles.

\figurePlaceholder{Figura o grafica del ejercicio. En un ejercicio real se puede incluir un SVG o PNG publico desde static/exercises/.}

\begin{exerciseParts}
  \item Identifica que magnitud cambia en el flujo magnetico.
  \item Calcula la magnitud fisica pedida en el enunciado.
  \item Justifica el sentido fisico o la interpretacion de la grafica, si procede.
\end{exerciseParts}

\sectionLabel{Pregunta de identificacion}

Antes de calcular, responde: que magnitud provoca la variacion del flujo magnetico?

\begin{itemize}
  \item Campo magnetico \(B\)
  \item Area efectiva \(A\)
  \item Angulo \(\theta\)
  \item Numero de espiras \(N\)
\end{itemize}

\end{document}
